\section{Del registro de trayectoria al operador de masa}
\label{vi:sec:operador-masa}

La construcción distingue dos operaciones. Primero, la dinámica produce
un registro de la trayectoria. Después, un funcional extrae coordenadas
enteras y las evalúa mediante las constantes internas ya obtenidas.
Conservar ese orden explica cómo la masa puede ser una magnitud de la
persistencia sin agotar su estructura.

\subsection{Lectura cronológica y eventos orientados}

Sea \(d_1,\ldots,d_{108}\) la lectura digital de una ruta enriquecida.
En la coordenatización declarada, el soporte de acción es \(\{1,3,5,7\}\), y
\[
 n_A=\#\{t:d_t\in\{1,3,5,7\}\}.
\]
Se conservan además los \(12\) eventos de los bloques de \(9\) pasos:
\[
 e_m=\sigma_m\#\{t:9m+1\le t\le9m+9,\ d_t\in\{2,4,6,8\}\},
 \quad 0\le m<12,
\]
donde \(\sigma=(+,+,+,-,-,-,+,+,+,-,-,-)\). El soporte de este
extractor es una parte explícita de la coordenatización de lectura; su función no
debe confundirse con el selector ternario de fase del TPK.

La distribución de los eventos es relevante incluso cuando sus sumas
coinciden. Para verla, introducimos
\[
 p_i=e_i+e_{i+6},\quad a_i=e_i-e_{i+6},\quad
 s_C=\sum_{i=0}^{5}p_i,\quad M_i=p_i-p_5\ (0\le i<5).
\]

\begin{proposition}[Reconstrucción integral de los eventos]
\label{vi:prop:eventos}
La transformación \((e_0,\ldots,e_{11})\mapsto(s_C,M_0,\ldots,M_4,a_0,\ldots,a_5)\)
es inyectiva sobre \(\mathbb Z^{12}\). Su imagen se caracteriza por
\[
 s_C-\sum_{i=0}^4M_i\equiv0\pmod6,
 \qquad p_i\equiv a_i\pmod2\quad(0\le i<6),
\]
donde \(p_5=(s_C-\sum M_i)/6\) y \(p_i=p_5+M_i\) para \(i<5\).
\end{proposition}
\begin{proof}
Sumar las seis expresiones de \(p_i\) da
\(s_C=6p_5+\sum_{i<5}M_i\). Esto recupera los \(p_i\), y después
\[
 e_i=\frac{p_i+a_i}{2},\qquad
 e_{i+6}=\frac{p_i-a_i}{2}.
\]
Las congruencias son necesarias para que esas cantidades sean enteras y
suficientes para que las fórmulas definan una preimagen integral. Sustituir
en la transformación directa recupera todas las coordenadas.
\end{proof}

La descomposición \(12=1+5+6\) adquiere así un significado operativo:
el total, las diferencias y las orientaciones recuperan los eventos.
La división por \(6\) pertenece a esa reconstrucción. Cuando el carácter
usa \(s_C/6\), \(s_C\) sigue siendo la suma entera; una coordenada ya
dividida no se divide una segunda vez.

\subsection{Memoria torsional y firma}

Para cada aparición de \(9\), la lectura conserva su predecesor.
Definimos, con índices cíclicos,
\[
 j_2(d)=\begin{cases}0,&d=9,\\1,&d\in\{2,4,6,8\},\\-1,&d\in\{1,3,5,7\},\end{cases}
 \quad z_t=\mathbf1_{\{d_t=9\}}j_2(d_{t-1}).
\]
Entonces
\[
 T_z=\sum_{t=1}^{108}z_t,\quad
 C_z=\sum_{t\in\{27,54,81,108\}}z_t,\quad
 k_{\Delta_4}=T_z-2C_z.
\]
La primera suma conserva la memoria acumulada y la segunda distingue los
cierres de corona. Eliminar los predecesores antes de evaluar estas
expresiones altera el objeto que se lee.

Con \(\tau_m=\operatorname{sgn}(e_m)\), las restantes coordenadas son
\[
 b_{120}=\sum_{a=0}^{3}\operatorname{sgn}(e_a+e_{a+4}+e_{a+8}),
 \qquad b_{\zeta}=\sum_{m=0}^{11}\tau_m\tau_{m+3},
\]
donde el segundo índice es módulo \(12\). Se obtiene la firma
\[
 L_{\mathrm{tick}}(\gamma)
 =\sigma(\gamma)=(n_A,s_C,k_{\Delta_4},b_{120},b_{\zeta})\in\mathbb Z^5.
\]
Esta aplicación es determinista en la coordenatización declarada. Los signos y las
correlaciones explican por qué la composición debe realizarse sobre el
registro antes de extraer los \(5\) enteros: la suma de firmas no
representa, en general, la composición de dos historias.

\subsubsection{Término de costura en la concatenación \(108\to1080\)}
\label{vi:subsubsec:costura-108-1080}

La prolongación de la lectura cronológica debe conservar el predecesor
efectivo de cada evento. Sean
\(d^{(b)}=(d^{(b)}_1,\ldots,d^{(b)}_{108})\), \(0\le b<10\), diez
palabras consecutivas, y definamos
\[
 a_b=d^{(b)}_1,\qquad q_b=d^{(b)}_{108},\qquad q_{-1}=q_9.
\]
Denotemos por \(T_z^{\sqcup}\) la suma de los diez lectores evaluados
por separado, usando \(q_b\) como predecesor cíclico de \(a_b\), y por
\(T_z^{\Vert}\) el lector de la palabra concatenada, que usa
\(q_{b-1}\) como predecesor efectivo de \(a_b\).

\begin{lemma}[Corrección exacta de costura]
\label{vi:lem:costura-108-1080}
Con \(\Delta T_z=T_z^{\Vert}-T_z^{\sqcup}\),
\begin{equation}
 \Delta T_z=
 \sum_{b=0}^{9}\mathbf1_{\{a_b=9\}}
 \bigl(j_2(q_{b-1})-j_2(q_b)\bigr).
 \label{vi:eq:delta-costura-tz}
\end{equation}
Si las coronas se sitúan en \(t\equiv0\pmod{27}\), entonces
\begin{equation}
 \Delta C_z=0,
 \qquad
 \Delta k_{\Delta_4}=\Delta T_z.
 \label{vi:eq:delta-costura-k}
\end{equation}
\end{lemma}

\begin{proof}
Fuera de los diez comienzos de bloque, cada evento conserva el mismo
predecesor en ambas evaluaciones. Si \(a_b\ne9\), el indicador que define
\(z_t\) anula también la posible diferencia de predecesor. Si \(a_b=9\),
el término separado es \(j_2(q_b)\) y el concatenado es
\(j_2(q_{b-1})\); su diferencia es el sumando de
\eqref{vi:eq:delta-costura-tz}. Al sumar los diez bloques se obtiene la
primera identidad. Los nuevos comienzos ocupan las posiciones
\(108b+1\equiv1\pmod{27}\), nunca una corona. Por tanto la costura no
modifica \(C_z\), y la definición
\(k_{\Delta_4}=T_z-2C_z\) da
\eqref{vi:eq:delta-costura-k}.
\end{proof}

La corrección alcanza realmente el valor \(-2\). Basta tomar un único
comienzo \(a_b=9\), con \(q_{b-1}=1\), \(q_b=2\), y exigir
\(a_r\ne9\) en los demás comienzos. Entonces
\(j_2(q_{b-1})-j_2(q_b)=-1-1=-2\). Esta palabra es un control del lector
de costura; no se promueve a ruta superviviente, partícula ni firma
física.

La longitud \(1080\) repite diez estaciones de fase, no diez copias de
una profundidad invariable. El transporte correcto es la composición
ordenada
\[
 \mathcal H_9\circ\mathcal H_8\circ\cdots\circ\mathcal H_0.
\]
Sólo una prueba adicional de estacionariedad permitiría sustituirla por
\(\mathcal H_{108}^{10}\). La fórmula de costura conserva precisamente
la memoria que esa sustitución podría borrar.


\subsection{Carácter y especialización interna}

Sea \(X=(X_A,X_C,X_4,X_{120},X_{270})\). El carácter formal es
\[
 \chi(\sigma;X)=\exp\!\left(n_AX_A+s_CX_C+k_{\Delta_4}X_4+
 b_{120}X_{120}+b_{\zeta}X_{270}\right).
\]
Para cada \(X\), satisface
\(\chi(\sigma+\sigma';X)=\chi(\sigma;X)\chi(\sigma';X)\)
porque el exponente es lineal en la firma. Es una ley sobre firmas; no
afirma que el extractor sea aditivo sobre rutas sin conservar la interfaz.

La especialización emplea las coordenadas internas construidas previamente:
\[
 X_A=A_{\mathrm{rad}},\quad X_C=C^*_{\mathrm{rad}}/6,
 \quad X_4=\Delta_4,\quad X_{120}=s_{120},\quad
 X_{270}=\zeta_{270}:=135\Delta_4^2.
\]
El carácter es adimensional y positivo. Su efecto es una homotecia sobre
la recta de masa. La forma exponencial determina cómo actúan los
coeficientes; su procedencia queda fijada en los registros y construcciones
anteriores.

\subsection{Jerarquía bicapa y conservación de las dos genealogías}

Escribimos \(x=A_{\mathrm{rad}}\), \(y=C^*_{\mathrm{rad}}\). En la
cámara \(x>y>0\), los canales son
\[
 q_+=e^{-x-y},\qquad q_-=e^{-x+y},\qquad0<q_+<q_-<1.
\]
Para una involución \(R\), con \(P_\pm=(I\pm R)/2\), el operador
de los dos canales es
\[
 T=q_+P_++q_-P_-.
\]
Sus momentos escalar y orientado se expresan como
\[
 c_n=q_-^n+q_+^n,\qquad s_n=q_-^n-q_+^n.
\]
Multiplicar estas expresiones prueba
\[
 c_n^2-s_n^2=4e^{-2nx},\quad
 c_{m+n}=\frac{c_mc_n+s_ms_n}{2},\quad
 s_{m+n}=\frac{s_mc_n+c_ms_n}{2}.
\]
Ambas sucesiones satisfacen
\[
 X_{n+2}=c_1X_{n+1}-e^{-2x}X_n,\qquad X\in\{c,s\}.
\]
Estas relaciones muestran que las alturas comparten un mismo operador.
No se define una colección distinta de momentos para cada especie.

\begin{proposition}[Semigrupo de alturas]
\[
 \langle90,120\rangle=\{90,120\}\cup\{30r:r\ge6\}.
\]
\end{proposition}
\begin{proof}
Tras dividir por \(30\), basta estudiar \(\langle3,4\rangle\).
Los enteros \(6,7,8\) son \(3+3,3+4,4+4\); sumar \(3\) produce
todos los enteros posteriores. Por debajo de \(6\), los únicos
elementos positivos son \(3,4\).
\end{proof}

Las alturas \(360=4\cdot90=3\cdot120\) ilustran la diferencia entre
igualdad de evaluación y memoria de composición. El cociente algebraico
identifica los dos exponentes; el registro conserva si el transporte
empleó cuatro etapas de un tipo o tres del otro.

La especialización refinada tiene exponente
\[
 \log\chi_{\mathrm{ref}}=n_Ax+\frac{s_C}{6}y+k_{\Delta_4}\Delta_4
 +b_{\zeta}\zeta_{270}+\sum_{h\in\langle90,120\rangle}N_hs_h,
 \qquad N_h=\eta_h+b_{120}\mathbf1_{\{h=120\}}.
\]
El par \((\eta_{120},b_{120})\) se conserva aunque la evaluación use
su suma. También se conservan separados \(s_{270}\) y
\(\zeta_{270}\): el primero es un momento bicapa y el segundo una
expresión cuadrática del refinamiento.

\begin{proposition}[Convergencia de la evaluación refinada]
Si \(\sum_h|N_hs_h|<\infty\), el carácter refinado define un escalar
positivo finito y sus evaluaciones truncadas convergen a él.
\end{proposition}
\begin{proof}
La hipótesis da convergencia absoluta del exponente. La continuidad de la
exponencial permite pasar al límite; la exponencial de un real finito es
estrictamente positiva.
\end{proof}

La convergencia evalúa una sucesión de coeficientes ya especificada.
La selección de esos coeficientes por la ruta es una operación distinta,
y debe acompañar a toda evaluación atribuida a una especie.

La notación \(\chi_{\mathrm{full}}\) conservada en algunas fórmulas
de evaluación de los sectores designa este mismo carácter
\(\chi_{\mathrm{ref}}\), con los mismos argumentos y coeficientes.
El cambio de símbolo no añade otra especialización ni otra torre.

\subsection{Lectores de desplazamiento y autocorrelación}

La misma ruta admite dos lecturas del retorno:
\[
 K_D=\left(D_{27}+D_{36},D_1,\frac{D_4-D_3}{2}\right),\qquad
 K_\Omega=(\Omega,54,P_6).
\]
El primer mapa conserva defectos de desplazamiento; el segundo, memoria de
retorno y autocorrelación. En el censo completo de \(324\) orientaciones
hay \(14\) perfiles del primero, \(9\) del segundo y \(40\) pares
conjuntos. Las \(18\) coincidencias tienen valor \((80,54,6)\).
No se infiere de ello la igualdad de ambos operadores en todo el atlas.

Para la ruta electrónica, los desplazamientos son
\[
 (D_1,D_3,D_4,D_{27},D_{36})=(54,56,68,48,32).
\]
La sustitución produce \(K_D=(80,54,6)\), coincidente con
\(K_\Omega\). Su exponente es
\[
 \kappa_e=80-54\alpha+6\Delta_4.
\]
La trayectoria electrónica se construye en
\eqref{vi:dep:trayectoria-electron}; sus lectores de desplazamiento y
autocorrelación se evalúan, respectivamente, en
\eqref{vi:dep:KD} y \eqref{vi:dep:KOmega}. La coincidencia
\(K_D=K_\Omega=(80,54,6)\) pertenece al teorema de cierre beta de la
subsección~\ref{vi:dep:beta}. Aquí esa excepción central se conserva y la
operación se extiende a las fibras del atlas.

\subsection{Operador sobre una fibra y sección dimensional}

Sea \(\mathscr H_M\) el espacio de Hilbert finito con base ortonormal
\((e_\gamma)\) etiquetada por las rutas de una multisección. Para
\(r\in\{D,\Omega\}\), el operador de retorno se define por
\(B_M^re_\gamma=\kappa_\gamma^re_\gamma\), con exponentes reales;
por tanto es autoadjunto. Sea \(M_M^{(0)}\) el operador basal positivo,
cuya conexión con el carácter se fija a continuación.
La razón adimensional de las
secciones de acción es \(R_{\mathrm{act}}>0\). Se define
\[
 M_M^{\beta,r}=R_{\mathrm{act}}^{B_M^r/2}
 M_M^{(0)}R_{\mathrm{act}}^{B_M^r/2}.
\]

\begin{proposition}[Positividad y restricción escalar]
El operador \(M_M^{\beta,r}\) es positivo. Si ambos operadores son
diagonales en la base de rutas, cada valor basal \(m_\gamma^{(0)}\)
se transforma en
\[
 m_\gamma^\beta=m_\gamma^{(0)}R_{\mathrm{act}}^{\kappa_\gamma}.
\]
\end{proposition}
\begin{proof}
Sea \(S=R_{\mathrm{act}}^{B_M^r/2}\), autoadjunto e invertible.
Para todo \(v\),
\(\langle v,SM_M^{(0)}Sv\rangle=\langle Sv,M_M^{(0)}Sv\rangle\ge0\).
En la base diagonal, los dos factores \(S\) multiplican cada valor
por \(R_{\mathrm{act}}^{\kappa_\gamma/2}\); su producto da la fórmula.
\end{proof}

La sección absoluta de masa se obtiene de la acción y el reloj:
\[
 \omega_0=\frac{2\pi}{108t_0},\qquad
m_{\mathrm{clk}}=\hbar_{\mathrm{ret}}\omega_0c_{\mathrm{int}}^{-2}.
\]
En la coordenatización normalizada respecto del electrón, el operador basal anterior es
\[
 M_M^{(0)}e_\gamma
 =m_e^{(0)}\,
   \chi_{\mathrm{ref}}(\sigma_\gamma-\sigma_e,
                       \eta_\gamma-\eta_e)e_\gamma.
\]
Aquí \(m_e^{(0)}\) es la realización de la coordenada electrónica
basal, con su prefactor, su refinamiento y su base dimensional.
Si \(E_e^{(0)}=\varepsilon_e^{(0)}\mathcal U_E\), entonces
\[
 m_e^{(0)}=\varepsilon_e^{(0)}\mathcal U_Ec_{\mathrm{int}}^{-2}
 =b_e\,m_{\mathrm{clk}},\qquad
 b_e=\varepsilon_e^{(0)}\frac{\mathcal U_E}{E_{\mathrm{clk}}}.
\]
En la coordenatización \(\mathcal U_E=E_{\mathrm{clk}}\) se obtiene
\(b_e=\varepsilon_e^{(0)}\); en otra coordenatización se conserva el cociente
de bases indicado. Si las coordenadas de torre ya son
relativas al electrón, se omite una segunda sustracción de \(\eta_e\).
Esta igualdad enlaza el operador con la firma de cada ruta y conserva
explícitamente la normalización electrónica. Una elección distinta de
sección de referencia debe transportar también esa normalización.
La escala de la constante reducida de Planck pertenece a esta construcción
dimensional. El carácter y el factor de retorno actúan después como
escalares. El electrón conserva su estructura: no se identifica con el
cuanto cronológico \(m_{\mathrm{clk}}\).

Para dos rutas en la misma coordenatización absoluta,
\[
 \frac{m_Y^\beta}{m_X^\beta}
 =\chi_{\mathrm{ref}}(\sigma_Y-\sigma_X,\eta_Y-\eta_X)
 R_{\mathrm{act}}^{\kappa_Y-\kappa_X}.
\]
La sección dimensional común cancela. La diferencia de los exponentes de
retorno no cancela salvo que sea nula. Esta identidad conserva la
distinción entre escala absoluta y estructura relativa de una colección.

\subsection{El sector de propagación nula}

La exponencial toma valores estrictamente positivos. La coordenatización nula se
construye antes de ese carácter, en la realización partida
\[
 \mathcal A_{\mathrm{sp}}=\mathbb R[R]/(R^2-1),\quad
 z=r_+P_++r_-P_-,\quad N(z)=r_+r_-.
\]
Las direcciones \(\mathbb RP_+\) y \(\mathbb RP_-\) tienen norma
nula. Una propagación transportada en una de ellas pertenece a ese sector;
la lectura material bicapa emplea el acoplamiento de ambas. De esta manera,
el fotón y el sector gluónico conservan su residencia en la teoría aunque
no constituyan una fila del carácter positivo. La realización de sus
representaciones se expone junto a los sectores correspondientes.
